\section{The \texttt{Source\_simple} McStas Component}
A circular neutron source with flat energy spectrum and arbitrary flux

\subsection*{Identification}
\begin{itemize}
  \item \textbf{Author:} Kim Lefmann
  \item \textbf{Origin:} Risoe
  \item \textbf{Date:} October 30, 1997
\end{itemize}

\subsection*{Description}
The routine is a circular neutron source, which aims at a square target centered at the beam (in order to improve MC-acceptance rate). The angular divergence is then given by the dimensions of the target. The neutron energy is uniformly distributed between lambda0-dlambda and lambda0+dlambda or between E0-dE and E0+dE. The flux unit is specified in n/cm2/s/st/energy unit (meV or \AA{}).

This component replaces Source\_flat, Source\_flat\_lambda, Source\_flux and Source\_flux\_lambda.

Example: Source\_simple(radius=0.1, dist=2, focus\_xw=.1, focus\_yh=.1, E0=14, dE=2)

\subsection*{Input parameters}
Parameters in \textbf{boldface} are required; the others are optional.

\begin{longtable}{p{0.22\textwidth}p{0.12\textwidth}p{0.46\textwidth}p{0.14\textwidth}}
\toprule
\textbf{Name} & \textbf{Unit} & \textbf{Description} & \textbf{Default} \\
\midrule
\endhead
radius & m & Radius of circle in (x,y,0) plane where neutrons are generated. & 0.1 \\
yheight & m & Height of rectangle in (x,y,0) plane where neutrons are generated. & 0 \\
xwidth & m & Width of rectangle in (x,y,0) plane where neutrons are generated. & 0 \\
dist & m & Distance to target along z axis. & 0 \\
focus\_xw & m & Width of target & .045 \\
focus\_yh & m & Height of target & .12 \\
E0 & meV & Mean energy of neutrons. & 0 \\
dE & meV & Energy half spread of neutrons (flat or gaussian sigma). & 0 \\
lambda0 & \AA{} & Mean wavelength of neutrons. & 0 \\
dlambda & \AA{} & Wavelength half spread of neutrons. & 0 \\
flux & 1/(s*cm**2*st*energy unit) & flux per energy unit, \AA{} or meV if flux=0, the source emits 1 in 4*PI whole space. & 1 \\
gauss & 1 & Gaussian (1) or Flat (0) energy/wavelength distribution & 0 \\
target\_index & 1 & relative index of component to focus at, e.g. next is +1 this is used to compute 'dist' automatically. & 1 \\
\bottomrule
\end{longtable}

\subsection*{Links}
\begin{itemize}
  \item Component source code found in file \texttt{Source\_simple.comp}.
\end{itemize}
\IfFileExists{sources/Source_simple_static.tex}{\input{sources/Source_simple_static.tex}}{}